\documentclass[11pt,oneside,openany]{book}

% 東吳大學法研所「憲法專題研究（一）」第 6 週（2026-10-19）導讀問答集：
% Jeremy Waldron, The Core of the Case Against Judicial Review, 115 Yale L.J. 1346 (2006)。
% 版式照 ~/Documents/NTU/1151/1151_Statistical_Deep_Learning/tex/1151_課程問答.tex。
% 宏包：Translation 倉 LaTeX/course-notes.sty；圖為 figures/*.tex 的 TikZ。
% 編譯：./build.sh（xelatex、biber、xelatex 兩次）。

\def\TransLaTeXDir{../../../../../1142/Translation/LaTeX}
\input{\TransLaTeXDir/course-notes.sty}
\addbibresource{refs.bib}

% 中文連續成段，行距取 1.5（LaTeX_排版設計基準）
\setstretch{1.5}
% 整理後的問題框改用中文標籤
\renewtcolorbox{qaquestion}{%
  enhanced, breakable, frame hidden, boxrule=0pt, arc=0pt,
  colback=white, borderline west={0.6pt}{0pt}{ruleink},
  left=9pt, right=8pt, top=2pt, bottom=2pt, before skip=2pt, after skip=10pt,
  fontupper=\setstretch{1.4}\setlength{\parskip}{0.25em},
  before upper={{\small\bfseries\color{dimink}問題}\par},
  before upper app={\setlist[enumerate]{label=\arabic*.,leftmargin=1.6em,itemsep=0.1em,topsep=0.2em,parsep=0pt}}}
% 導讀人自己的推論，與原文所述分開
\newtcolorbox{readernote}{%
  enhanced, breakable, frame hidden, boxrule=0pt, arc=0pt,
  colback=figGraybg, left=9pt, right=8pt, top=4pt, bottom=4pt,
  before skip=6pt, after skip=8pt,
  fontupper=\small\setstretch{1.4}\setlength{\parskip}{0.2em},
  before upper={{\small\bfseries\color{figGray}導讀人按}\par}}
% 刊本頁
\newcommand{\pg}[1]{（#1）}
\newcommand{\en}[1]{\textit{#1}}
% 章標改中文：第 N 部
\titleformat{\chapter}[display]
  {\normalfont\color{caseink}}
  {\small\color{dimink}第 \thechapter{} 部}
  {4pt}
  {\LARGE\bfseries}
% 西文長句與不可斷的頁碼框，給行調整一點餘地
\setlength{\emergencystretch}{2em}
% 圖與表的標籤
\captionsetup[figure]{name=圖}
\captionsetup[table]{name=表}
% 圖的共用樣式（diagram-design skill）：灰階細框為底，強調色 figE 一圖只用一處
\tikzset{
  qa/.style={font=\footnotesize,every node/.append style={execute at begin node={\setstretch{1.2}}}},
  nbox/.style={draw=figGray,line width=0.4pt,fill=white,rounded corners=1.5pt,inner xsep=3pt,inner ysep=2.6pt,align=left},
  gbox/.style={nbox,fill=figGraybg},
  kbox/.style={nbox,draw=figE,fill=figEbg,line width=0.7pt},
  lbl/.style={font=\scriptsize,text=dimink,inner sep=1pt},
  hdr/.style={font=\footnotesize,text=caseink,inner sep=1pt},
  arr/.style={figGray,line width=0.6pt,-{Stealth[length=4pt,width=3.2pt]},shorten >=1pt,shorten <=1pt},
  dsh/.style={arr,dash pattern=on 2.4pt off 1.6pt},
}
% 表格窄欄齊左
\newcolumntype{L}[1]{>{\raggedright\arraybackslash}p{#1}}

\newcommand{\documentversion}{v0.1（2026-10-18）}

\begin{document}

\frontmatter
\coursetitle{反對違憲審查的核心論據}{Waldron（2006）導讀問答集}%
  {東吳大學法律學系研究所　憲法專題研究（一）\\115 學年度第 1 學期　第 6 週（2026 年 10 月 19 日）}%
  {\documentversion}
\tableofcontents

\chapter*{凡例}
\addcontentsline{toc}{chapter}{凡例}
\begin{itemize}
\item 本書所讀為 Jeremy Waldron, The Core of the Case Against Judicial Review, 115 Yale L.J. 1346--1406 (2006)。課綱作 The Core Case against Judicial Review，刊本篇名多 of the 二字，書目依刊本。
\item 正文括號內的四位數字為刊本頁，例如\pg{1387}。原文 PDF 頁加 1344 即刊本頁。
\item 譯名照譯本《反對違憲審查的核心論據》（Translation 倉 JudicialReview）的譯語表；引文的中譯取自該譯本。
\item 各則依本文自己的論證順序排列，前一部回答導讀開場的四個問題，後三部放回本課的課綱、2006 年以後的發展與討論題。
\item 灰底「導讀人按」為導讀人的推論與對照，原文未言及；其餘各句均可指回原文頁碼或所引文獻。
\end{itemize}

\mainmatter
% 第一部：作者、寫作脈絡、對話對象與本文的新意。
% 出處：刊本作者註（1346）、導論（1348–1353）；課程倉 notes/docs/1151/seminar/W06_1019_導讀問答.md 2026-10-06 兩則與其查核報告。
\chapter{作者、寫作脈絡與對話對象}

\qaentry{第一部之一}{作者 Jeremy Waldron}

\begin{qaquestion}
Waldron 撰文時任職何處？學術訓練與經歷為何？\\
這些背景在本文留下哪些痕跡？
\end{qaquestion}

\subsection{任職與經歷}

刊本的作者註寫明，Waldron 撰文時是哥倫比亞大學法學院的 University Professor，2006 年 7 月起轉任紐約大學法學院教授\pg{1346}。紐約大學法學院 2006 年的院刊介紹他時寫到：他在紐西蘭南部的 Invercargill 長大，在 Otago 大學取得哲學與法律學位，1978 年赴牛津攻讀研究所，1980 至 1982 年任牛津 Lincoln College 研究員，其後任教於愛丁堡大學、柏克萊、普林斯頓與哥倫比亞；同一篇介紹以 \en{taking strong, sometimes controversial stands on judicial review (should be weakened)} 概括他在違憲審查上的立場。\footcite{nyumag2006}

\subsection{紐西蘭背景在本文的痕跡}

紐西蘭之名在正文與註中出現於七頁（1349、1356、1361、1362、1365、1370、1394），用途分成兩組。

\begin{itemize}
\item 正面的例子。英國國會 1960 年代就墮胎、同性性行為、死刑所作的辯論，與加拿大、澳洲、紐西蘭等地的類似辯論，被用來反駁「議員無法負責任地處理這類問題」之說\pg{1349}；1990 年權利法案法是「更弱的形式」違憲審查的實例\pg{1356}；1993 年國會立法把選制改為比例代表制，說明立法機關能自行改革選舉制度\pg{1362~註~48}；權利法案法第 7 條要求檢察總長在法律案提出後指出其中與權利法案不符的條文，作者說這類安排是他所能接受的\pg{1370~註~63}。
\item 反面的例子。註 47 引他自己 2005 年在 New Zealand Law Journal 的短文，表示紐西蘭的一院制可能已使該國落在本文論證的範圍之外\pg{1361}；第柒部分註 136 再引這篇短文，作為論證不成立的一例\pg{1402}。
\end{itemize}

\begin{readernote}
作者用四項假設檢驗自己的國家，結論是紐西蘭未必合格。導讀時可以把這兩組並列：他所主張的制度在紐西蘭有實例，而他也承認紐西蘭的國會可能不符合第一項假設。
\end{readernote}

\qaentry{第一部之二}{作者為何在 2006 年再寫一篇反對違憲審查的文章}

\begin{qaquestion}
導論說雙方陣線已定，作者本人也寫過不少，為何再寫一篇？他自己給了哪些理由？
\end{qaquestion}

導論自問此事\pg{1350--1351}，答以四項理由。

\begin{enumerate}
\item 找出獨立於歷史表現與個案結果的核心論證。他寫到，Tushnet 與 Kramer 的近著把對違憲審查的理論批評，與它的歷史起源、以及兩人對一部較少司法化的美國憲法的構想交纏在一起，並表明這不是對兩人的批評。Michelman 為 Kramer 一書所寫的推薦語說，該書的歷史 \en{puts flesh on the bones of debates over judicial review}；Waldron 接著說，他要 \en{take off some of the flesh and boil down the normative argument to its bare bones}\pg{1351}。
\item 讓論辯脫離對個別判決的好惡。他引 Charles Black 的話，人們對違憲審查的反對往往是 \en{a sometime thing}，並舉墮胎權論辯為例：支持選擇權的一方擔心動搖 Roe v.\ Wade，連反對違憲審查的論證都不願考慮\pg{1351--1352}。
\item 讓論證適用於美國以外的國家。英國人討論本國法官有限的審查權時，不關心 1805 年聯邦黨人與共和黨人的爭論，也不關心 Brown 的遺產\pg{1352}。
\item 對自己先前的表述不滿意。他寫道，自己過去的討論把違憲審查與法理學、政治哲學的其他問題交纏在一起，未曾清楚陳述基本的反對理由，也未充分回應對 Law and Disagreement 等著作的批評\pg{1352--1353}。
\end{enumerate}

摘要把全文的主張歸成兩項：其一，沒有理由認為權利在違憲審查之下比在民主立法機關之下受到更好的保護；其二，撇開所產生的結果，違憲審查在民主上欠缺正當性。摘要自述第二項為人熟知，第一項較少人提出\pg{1346}。導論末段把同一主張寫成兩條戰線：違憲審查以先例、文本與解釋的枝節問題分散公民對真正爭點的注意；它讓少數非民選、不須負責的法官以多數決作終局決定，剝奪一般公民的參政權，撇開代表與政治平等的原則\pg{1353}。

\qaentry{第一部之三}{寫作時的論辯脈絡}

\begin{qaquestion}
本文發表時，美國與其他普通法國家的違憲審查論辯處於什麼狀態？本文在三次發表場合面對的是哪些人？
\begin{enumerate}
\item 美國的論辯。
\item 大英國協的新模式。
\item 發表場合與評論人。
\end{enumerate}
\end{qaquestion}

\begin{qafigure}{本文前後的相關事件。年份與刊本頁的出處見本則各小節。\label{fig:timeline}}
% 本文前後的相關事件：蛇行時間軸，每列四則（改自 ~/.claude/skills/diagram-design/templates/蛇行時間軸.tex）。
% 灰階為底，本文一則用強調色。事件卡第一行小字為類別。
\begin{tikzpicture}[qa,
    track/.style={line width=1pt,color=ruleink,line cap=round},
    dot/.style={circle,fill=#1,draw=white,line width=1pt,inner sep=0pt,minimum size=6.5pt},
    dt/.style={font=\footnotesize\bfseries,text=#1,inner sep=1pt},
    card/.style={nbox,text width=29mm,font=\scriptsize,text=caseink,anchor=north,
      execute at begin node={\raggedright\setstretch{1.2}}},
    keycard/.style={card,kbox}]
  \def\xa{1.9}\def\dx{3.4}\def\dy{3.0}\def\R{1.5}
  \pgfmathsetmacro{\xl}{\xa-0.7}\pgfmathsetmacro{\xr}{\xa+3*\dx+0.7}
  \draw[track] (\xl,0) -- (\xr,0) arc[start angle=90,end angle=-90,radius=\R]
    -- (\xl,-\dy) arc[start angle=90,end angle=270,radius=\R]
    -- (\xr,-2*\dy) arc[start angle=90,end angle=-90,radius=\R]
    -- ({\xa+1.6*\dx},-3*\dy);
  \draw[track,-{Stealth[length=2.4mm,width=2.2mm]}] ({\xa+1.6*\dx},-3*\dy) -- ({\xa+1.4*\dx},-3*\dy);
  \foreach \r/\c/\k/\d/\t in {
      0/0/n/1938/{\textcolor{dimink}{判決}\\Carolene Products 案腳註四：孤立而隔絕的少數},
      0/1/n/1962/{\textcolor{dimink}{著作}\\Bickel, The Least Dangerous Branch：對抗多數的難題},
      0/2/n/1982/{\textcolor{dimink}{立法}\\加拿大權利與自由憲章第 33 條},
      0/3/n/1990/{\textcolor{dimink}{立法}\\紐西蘭權利法案法},
      1/0/n/1998/{\textcolor{dimink}{立法}\\英國人權法：不相容宣告},
      1/1/n/1999/{\textcolor{dimink}{著作}\\Waldron, Law and Disagreement；Tushnet, Taking the Constitution Away from the Courts},
      1/2/n/2003/{\textcolor{dimink}{判決}\\Goodridge 案：麻州同性婚姻},
      1/3/n/2004/{\textcolor{dimink}{著作}\\Kramer, The People Themselves},
      2/0/k/2006/{\textcolor{dimink}{本文}\\115 Yale L.J. 1346},
      2/1/n/2008/{\textcolor{dimink}{著作}\\Fallon 以回應本文為題},
      2/2/n/2014/{\textcolor{dimink}{著作}\\Waldron, Five to Four：法院內部的多數決},
      2/3/n/2015/{\textcolor{dimink}{判決}\\Obergefell 案，五比四},
      3/0/n/2022/{\textcolor{dimink}{判決}\\Dobbs 案：推翻 Roe 與 Casey},
      3/1/n/2022/{\textcolor{dimink}{立法}\\紐西蘭權利法案法增訂第 7A 條}}{
    \pgfmathtruncatemacro{\odd}{mod(\r,2)}
    \pgfmathsetmacro{\x}{\odd==0 ? \xa+\c*\dx : \xa+(3-\c)*\dx}
    \pgfmathsetmacro{\y}{-\r*\dy}
    \if k\k
      \node[dot=figE] at (\x,\y) {};
      \node[dt=figE,anchor=south] at (\x,\y+0.16) {\d};
      \node[keycard] at (\x,\y-0.22) {\t};
    \else
      \node[dot=caseink] at (\x,\y) {};
      \node[dt=caseink,anchor=south] at (\x,\y+0.16) {\d};
      \node[card] at (\x,\y-0.22) {\t};
    \fi
  }
\end{tikzpicture}

\end{qafigure}

\subsection{美國的論辯}

導論描述了 2006 年前後的美國論辯\pg{1350}。長期以來，支持違憲審查的是自由派，反對的是不滿自由派法院所支持之權利的保守派；Rehnquist 法院廢棄若干自由派立法之後（註~12 舉 United States v.\ Morrison、United States v.\ Lopez），自由派的反對也增加了。註 11 列 Kramer 的 The People Themselves（2004）與 Tushnet 的 Taking the Constitution Away from the Courts（1999）為近年攻擊違憲審查之書，Dworkin、Eisgruber、Sager 則為違憲審查辯護（註~\mbox{11、13}）。Marbury v.\ Madison 兩百週年與 Brown v.\ Board of Education 五十週年，又帶出一批關於違憲審查起源與成就的討論。

\subsection{大英國協的新模式}

第壹部分界定強式與弱式違憲審查時，除美國與部分歐洲法院之外，另舉大英國協的三例：英國 1998 年人權法的不相容宣告為弱式\pg{1355}，紐西蘭 1990 年權利法案法只要求法院儘可能作合於權利的解釋，為更弱的形式\pg{1356}，加拿大權利與自由憲章第 33 條的「儘管條款」雖使立法機關得排除審查，作者仍計為強式\pg{1356--1357}。本文的攻擊對象限於強式審查，前兩種制度提供了作者所能接受的對照（見第三部）；加拿大屬於攻擊的標的。

\subsection{發表場合與評論人}

作者註列出本文的三次發表：倫敦大學學院的法律與社會哲學研討會、耶路撒冷希伯來大學的法學院工作坊、哈佛法學院的憲法研討會，Dworkin、Ruth Gavison、Seana Shiffrin 擔任正式評論人；致謝名單中還有 Raz、Fallon、Michelman 等人\pg{1346}。正文所回應的論者，有幾位就在這份名單上。

\qaentry{第一部之四}{本文與作者先前著作的關係}

\begin{qaquestion}
本文是 Waldron 違憲審查理論的首次陳述，還是重述？與 Law and Disagreement 等書是什麼關係？
\end{qaquestion}

註 14 列出他此前討論違憲審查的著作，包括 Law and Disagreement（1999）第 10 至 17 頁與第 211 至 312 頁、Deliberation, Disagreement, and Voting（1999）、Judicial Power and Popular Sovereignty（2002）；註 19 另列他從權利理論、公民共和主義、盧梭與邊沁兩種民主觀、歐陸人民主權論等角度處理同一問題的文章\pg{1351--1353}。\footcite{waldron1999lad} 第肆部分註 89 引 The Dignity of Legislation（1999），說明他刻意以樂觀的立法機關圖像對應法院支持者的樂觀圖像\pg{1379}。\footcite{waldron1999dignity}

依作者自述\pg{1352--1353}，本文重述他已有的立場，把論證從上列著作與政治哲學的其他問題中抽出，寫成可以單獨檢驗的形式。四項假設、結果面與程序面兩類理由的區分、核心情形與非核心情形的劃分，都在本文寫成可以單獨檢驗的形式；註 42 說明四項假設改寫自作者 2004 年 Some Models of Dialogue Between Judges and Legislators 一文。

\qaentry{第一部之五}{本文和誰對話}

\begin{qaquestion}
本文引用、回應了哪些論者？依立場可分幾群？作者對各說如何回應？
\end{qaquestion}

\begin{qafigure}{本文的對話對象，依作者對各人的態度分三組。長條為本文正文與註中提及該姓氏的次數，取自課程倉 \texttt{literature-mentions.tsv}。\label{fig:interlocutors}}
% 本文的對話對象：三組，依位置分組，灰階；長條長度＝本文提及次數（literature-mentions.tsv）。
\begin{tikzpicture}[qa]
  \def\u{0.12}% 每次提及 1.2mm
  \newcommand{\who}[4]{% x, y, 名, 次數（0 表示未計）
    \node[anchor=east,font=\footnotesize] at (#1,#2) {#3};
    \ifnum#4>0
      \fill[figGray] (#1+0.1,#2-0.07) rectangle ++({#4*\u},0.14);
      \node[lbl,anchor=west] at ({#1+0.15+#4*\u},#2) {#4};
    \fi}
  \node[hdr,anchor=west] at (0,0.6) {所反駁的辯護者};
  \who{2.2}{0}{Dworkin}{28}
  \who{2.2}{-0.45}{Rawls}{17}
  \who{2.2}{-0.9}{Raz}{11}
  \who{2.2}{-1.35}{Eisgruber}{7}
  \who{2.2}{-1.8}{Moore}{0}
  \who{2.2}{-2.25}{Christiano}{0}
  \node[hdr,anchor=west] at (8.4,0.6) {同樣批評違憲審查};
  \who{10.6}{0}{Tushnet}{10}
  \who{10.6}{-0.45}{Kramer}{8}
  \node[hdr,anchor=west] at (8.4,-0.95) {借用其概念或論證工具};
  \who{10.6}{-1.4}{Michelman}{6}
  \who{10.6}{-1.85}{Bickel}{3}
  \node[lbl,anchor=north west,text width=60mm,align=left] at (8.4,-2.15) {Kavanagh、Fuller、Wollheim、Condorcet、Elster、Scalia（Casey 案不同意見）};
\end{tikzpicture}

\end{qafigure}

\subsection{作者所反駁的辯護者}

\begin{longtable}{@{}L{0.17\linewidth}L{0.36\linewidth}L{0.39\linewidth}@{}}
\toprule
論者 & 主張 & 作者的回應 \\
\midrule
\endhead
Dworkin & 法院是「原則的論壇」\pg{1385--1386}；參與的價值在於確認公民平等成員地位的後果\pg{1375}；尊重權利是任何決策制度具正當性的背景條件\pg{1398--1400} & 法院的推理以解釋與教義為主；Dworkin 嚴重低估了參與權的觀念；「彼此尊重權利」有兩種讀法，作者採「認真對待並真誠設法確認」一讀，在核心情形下已由第三項假設滿足 \\
Raz & 結果面理由是唯一值得考慮的理由\pg{1374}；立法機關受宗派利益驅使\pg{1376--1378} & 前者被稱為教條主義；同樣的宗派壓力也見於法院（Korematsu 等案），立法機關若普遍如此，即與第三項假設不相容，屬非核心情形 \\
Moore & 法官每天面對具體個人的處境，較能得到道德洞見\pg{1379} & 案件上訴到最高審級時，原告的個人處境已幾乎不見，爭的是抽象的權利問題\pg{1380} \\
Eisgruber & 法官經提名同意而有民主資格\pg{1394~註~125} & 法官的民主資格遠不及議員 \\
Rawls & 良序社會亦宜設立違憲審查\pg{1366~註~53} & 作者的目標之一是證明 Rawls 對良序社會與違憲審查的看法有誤；同時借用 Rawls 的「判斷的負擔」作為第四項假設的依據\pg{1368} \\
\bottomrule
\end{longtable}

\subsection{同樣批評違憲審查的近著}

Tushnet 與 Kramer 被列為近年攻擊違憲審查的著作\pg{1350~註~11}。作者表明不批評兩書，只把其中的歷史與制度構想抽去，留下規範論證\pg{1351}。

\subsection{借用其概念或論證工具者}

\begin{itemize}
\item Bickel 的「對抗多數的難題」與「脫軌制度」，作為既有反對論的代表\pg{1349}。\footcite[16--18]{bickel1962}
\item Michelman 所提的問法（何種方法最可能觸及權利真相、同時充分尊重受影響者平等發言的主張），作者將其前後倒置成第二種問法，用以設定比較兩類理由的問題\pg{1375~註~76}。
\item Kavanagh：結果面理由不必預設權利內容，只問何種程序較可能觸及權利的真相\pg{1373--1374}。
\item Fuller 與 Hart、Sacks 的法律程序學派，作為第肆部分檢驗機構能力的起點\pg{1376}。
\item Wollheim 的民主理論弔詭，作者改說成一切政治理論都有的弔詭，用以說明任何決定程序都會出錯\pg{1372}；Christiano 的程序無窮後退則是作者所回應的反對意見，作者指 Christiano 未證明其為惡性後退\pg{1371~註~66}。
\item Condorcet 陪審團定理，作者指出它反而有利於人數較多的立法機關\pg{1392}。
\item Elster 的 Ulysses 之喻，見於五項調和論證中的預先承諾說\pg{1393--1394}。
\item Scalia 在 Casey 案的不同意見，作者據以說明法院的正當性問題首當其衝\pg{1390--1391}。
\end{itemize}

\qaentry{第一部之六}{本文的新意}

\begin{qaquestion}
與 Bickel 以來的反對論相比，本文提出了哪些新的論證工具？各在何處？
\end{qaquestion}

摘要自述，主張違憲審查不能更好地保護權利的那一項較少人提出\pg{1346}。依作者在本文所用的工具，可整理成下表。

\begin{longtable}{@{}L{0.24\linewidth}L{0.58\linewidth}L{0.10\linewidth}@{}}
\toprule
工具 & 內容 & 刊本頁 \\
\midrule
\endhead
核心論據的寫法 & 抽離歷史表現與個案結果，只留規範論證 & 1351 \\
附條件的論證 & 論證以四項假設為前提；假設全部成立為核心情形，任一不成立屬非核心情形，論證可能不成立 & 1359--1360 \\
以權利分歧為出發點 & 社會成員認真對待權利，而就有哪些權利、其內容為何存在善意的實質分歧；與道德相對主義無關 & 1366--1369 \\
兩類理由 & 結果面問何者較可能得出正確結果；程序面問何人有權參與決定，與結果無涉 & 1372--1373 \\
C\textsubscript{n} 的兩個問題 & 向不同意決定的公民說明：為何由這些人決定、為何採此決定程序 & 1387 \\
法院內部的多數決 & 大法官之間的簡單多數決欠缺公平性的論證 & 1391--1392 \\
兩種多數與少數 & 決定上的少數與議題上的少數重合時，始構成多數暴政 & 1397--1398 \\
\bottomrule
\end{longtable}

\begin{readernote}
前三項決定了論證的形狀：先限定爭點，再限定社會，最後只比較兩個機關。後四項是比較時用的尺。導讀時可以先講前三項，聽者才知道後面的比較在什麼條件下成立。
\end{readernote}

% 第二部：篇名與全文架構。出處：課程倉 1151/seminar/my/導讀_W06_Waldron_細目次與大綱.md 第二至四節；翻譯倉附錄圖「全文的論證結構」。
\chapter{篇名與全文架構}

\qaentry{第二部之一}{篇名的 the core of the case 與正文的 core case}

\begin{qaquestion}
篇名的 the core of the case 指什麼？正文另有 core case 一詞，兩者是否同義？
\end{qaquestion}

篇名的 the core of the case 指反對違憲審查的論據中，抽離各國歷史與個案結果之後所餘的部分\pg{1351}，譯本譯作「反對違憲審查的核心論據」。正文的 core case 指四項假設全部成立的社會\pg{1359--1360}，譯作「核心情形」，與第柒部分的 non-core cases（非核心情形）成對。第參部分末段把兩者接在一起：結果面的理由至多未有定論，程序面的理由幾乎全部不利於違憲審查，立法機關一方在兩種問法下都勝出，作者說 \en{that will be the core of the case against judicial review}\pg{1375--1376}。

\qaentry{第二部之二}{七個部分的排序}

\begin{qaquestion}
本文的七個部分各回答什麼問題？排序本身構成怎樣的論證？
\end{qaquestion}

\begin{qafigure}{全文的論證結構（1353--1406）。改繪自譯本附錄〈全文的論證結構〉。\label{fig:structure}}
% 全文的論證結構，改繪自翻譯倉 JudicialReview/Werke/Teile/93_Argument_Figures.tex。灰階為底，結論用強調色。
\begin{tikzpicture}[qa,node distance=4.5mm]
  \newcommand{\p}[1]{\\[-1pt]{\scriptsize\color{dimink}#1}}
  \node[gbox,text width=30mm] (a1) at (0,0) {一、民主制度\p{1361--1362}};
  \node[gbox,text width=30mm,right=2mm of a1] (a2) {二、司法制度\p{1363--1364}};
  \node[gbox,text width=30mm,right=2mm of a2] (a3) {三、對權利的承諾\p{1364--1366}};
  \node[gbox,text width=30mm,right=2mm of a3] (a4) {四、關於權利的分歧\p{1366--1369}};
  \node[hdr,anchor=south west] at ([yshift=1.2mm]a1.north west) {貳、四項假設};
  \coordinate (bar) at ([yshift=-3.5mm]a1.south);
  \node[nbox,text width=64mm,anchor=north] (core) at ([yshift=-4mm,xshift=24mm]bar) {核心情形：四項假設均成立\p{1359--1360}};
  \node[nbox,text width=80mm,below=of core] (form) {參、論證的形式：權利分歧須經決定程序解決；評價程序的理由分為結果面與程序面\p{1369--1376}};
  \node[nbox,text width=38mm,below=of form.south west,anchor=north west] (out) {肆、結果面的理由：個案取向、文本取向、說明理由；至多未有定論\p{1376--1386}};
  \node[nbox,text width=38mm,below=of form.south east,anchor=north east] (proc) {伍、程序面的理由：由誰決定、以何種程序決定；大多有利於立法機關\p{1386--1395}};
  \node[nbox,text width=80mm,anchor=north] (tyr) at ([yshift=-4.5mm]out.south -| form) {陸、多數暴政：在第三、四項假設之下，不應預期議題上的少數與決定上的少數重合\p{1395--1401}};
  \node[kbox,text width=80mm,below=of tyr] (con) {結論：權利分歧應以立法制度解決，以權利為本的違憲審查不適合這樣的社會\p{1360；1406}};
  \node[nbox,text width=31mm,anchor=north west] (nc) at ([xshift=5mm]form.north east) {柒、非核心情形：民主制度失能、成員甚少關心他人權利、孤立而隔絕的少數；反對論不成立，違憲審查也不因此成立\p{1401--1406}};
  \foreach \x in {a1,a2,a3,a4} \draw[figGray,line width=0.5pt] (\x.south) -- (\x.south |- bar);
  \draw[figGray,line width=0.5pt] (a1.south |- bar) -- (a4.south |- bar);
  \draw[arr] (bar -| core) -- (core);
  \draw[arr] (core) -- (form);
  \draw[arr] (form.south -| out) -- (out);
  \draw[arr] (form.south -| proc) -- (proc);
  \draw[arr] (out.south) -- (out.south |- tyr.north);
  \draw[arr] (proc.south) -- (proc.south |- tyr.north);
  \draw[arr] (tyr) -- (con);
  \draw[dsh] (bar -| nc) -- node[lbl,fill=white,right] {假設不成立} (nc.north);
  \draw[dsh] (tyr.east) -| node[lbl,fill=white,pos=0.75,right] {兩種少數普遍重合} (nc.south);
\end{tikzpicture}

\end{qafigure}

導論末段已列出同一順序\pg{1353}。依各部分回答的問題，排序如下。

\begin{longtable}{@{}L{0.25\linewidth}L{0.30\linewidth}L{0.37\linewidth}@{}}
\toprule
部分 & 要回答的問題 & 作用 \\
\midrule
\endhead
壹、違憲審查的界定\newline\textcolor{dimink}{\small 1353--1359} & 攻擊的是哪一種違憲審查 & 標的：限於對立法的強式審查 \\
貳、四項假設\newline\textcolor{dimink}{\small 1359--1369} & 在什麼樣的社會裡討論 & 前提：四項假設界定核心情形 \\
參、論證的形式\newline\textcolor{dimink}{\small 1369--1376} & 以什麼標準選擇決定程序 & 判準：結果面與程序面兩類理由 \\
肆、結果面的理由\newline\textcolor{dimink}{\small 1376--1386} & 法院是否較可能得出正確的權利決定 & 檢驗一：法院被歸予的三項優點 \\
伍、程序面的理由\newline\textcolor{dimink}{\small 1386--1395} & 撇開結果，誰有權作終局決定 & 檢驗二：向不同意者說明的兩個問題 \\
陸、多數暴政\newline\textcolor{dimink}{\small 1395--1401} & 多數暴政之說能否支持違憲審查 & 對最常見反駁的回應 \\
柒、非核心情形\newline\textcolor{dimink}{\small 1401--1406} & 假設不成立時如何 & 前提不成立時論證的效力 \\
\bottomrule
\end{longtable}

第貳部分的四項假設各有後文取用之處：B 節的法院特徵在第肆部分 A 節與第伍部分取用，C 節在第陸、柒部分取用，D 節在第參部分與第肆部分 B 節取用，A 節在第伍、柒部分取用。非核心情形在第柒部分之前已多次預告\pg{1360、1366、1377、1386、1398}。

\qaentry{第二部之三}{小標未顯示的內容}

\begin{qaquestion}
刊本只有兩層共 16 條小標。導讀時有哪些內容在小標上看不出來，應向聽者指出？
\end{qaquestion}

\begin{enumerate}
\item 第壹部分只有一個小標，內容卻有四項區分：強式與弱式、以權利為對象的審查與一般憲法審查、事後的個案審查與事前的抽象審查、普通法院與專設的憲法法院；後三項各帶作者的評論\pg{1354--1359}。
\item 第伍部分後半回應五項調和違憲審查與民主的論證，小標未列\pg{1393--1395}。
\item 第陸部分末尾回應 Dworkin 的「背景條件」說與 nemo iudex in sua causa（任何人不得為自己案件之法官）\pg{1398--1401}。
\end{enumerate}

% 第三部：攻擊的標的（第壹部分，1353–1359）。出處：課程倉 W06_1019_導讀問答.md「強式與弱式」一則（已查核）。
\chapter{攻擊的標的：違憲審查的界定}

\qaentry{第三部之一}{強式與弱式的區分標準}

\begin{qaquestion}
本文攻擊的是哪一種違憲審查？強式與弱式以什麼為區分標準？
\end{qaquestion}

第壹部分開頭先限定範圍：本文處理法院對立法的審查，對行政行為的司法審查依據的是法治原則，不在本文範圍\pg{1353--1354}。在對立法的審查中，作者認為最重要的區分是強式與弱式\pg{1354}。強式的定義是：

\begin{quote}
In a system of strong judicial review, courts have the authority to decline to apply a statute in a particular case (even though the statute on its own terms plainly applies in that case) or to modify the effect of a statute to make its application conform with individual rights (in ways that the statute itself does not envisage).\pg{1354}
\end{quote}

法院另得經由遵循先例原則與爭點效，使該法律實際上成為具文。弱式的定義是：

\begin{quote}
In a system of weak judicial review, by contrast, courts may scrutinize legislation for its conformity to individual rights but they may not decline to apply it (or moderate its application) simply because rights would otherwise be violated.\pg{1355}
\end{quote}

區分的標準因此是法院在個案中得否拒絕適用法律，或修正其效果。註 22 另指出，強弱之分與司法至上（judicial supremacy）屬不同問題\pg{1354}。

\qaentry{第三部之二}{各國制度的歸類}

\begin{qaquestion}
本文如何把美國、加拿大、英國、紐西蘭與部分歐洲法院排在強弱之間？加拿大有儘管條款，為何仍算強式？
\end{qaquestion}

\begin{qafigure}{違憲審查的五種形式與原文所舉之例（1354--1357）。\label{fig:spectrum}}
% 違憲審查的五種形式，由強至弱（1354–1357）。灰階；本文攻擊的標的以強調色括號標出。
\begin{tikzpicture}[qa]
  \foreach \i/\n/\d/\e in {
      0/更強/{將法律自法典中剔除}/{部分歐洲法院\\1354--1355 註 23},
      1/強式/{個案中拒絕適用或修正其效果；經先例使之成為具文}/{美國\\1355},
      2/居間，計為強式/{拒絕適用；立法機關得以儘管條款排除}/{加拿大憲章第 33 條\\1356--1357},
      3/弱式/{不得拒絕適用；得為不相容宣告}/{英國 1998 年人權法\\1355},
      4/更弱/{只得作合於權利的解釋}/{紐西蘭 1990 年權利法案法\\1356}}{
    \node[nbox,text width=25mm,minimum height=21mm,anchor=north west] (b\i) at (\i*2.95,0) {\n\\[2pt]{\scriptsize \d}};
    \node[lbl,text width=26mm,align=left,anchor=north west] at (\i*2.95,-2.35) {\e};
  }
  \draw[figGray,line width=0.6pt,-{Stealth[length=4.5pt]}] (0,0.45) -- node[lbl,above] {法院的權限由強至弱} (14.6,0.45);
  \draw[figE,line width=0.9pt] (0.05,-3.25) -- (8.5,-3.25);
  \node[font=\scriptsize,text=figE,anchor=north west] at (0.05,-3.32) {本文攻擊的標的};
\end{tikzpicture}

\end{qafigure}

美國法院得拒絕適用，依原文似不能將法律剔除，但實際效果與剔除相去不遠\pg{1355}。註 24 討論美國是否屬於更強的形式，稱此事並無明確答案：Dickerson v.\ United States 案 Scalia 大法官不同意見書的末句，似乎顯示被最高法院認定違憲的法律仍可供法院援用；McCorvey v.\ Hill 案中第五巡迴上訴法院認定德州墮胎法已遭默示廢止，作者細讀後指出，默示廢止出於 Roe 之後德州的墮胎立法，並非 Roe 判決本身\pg{1355}。

加拿大權利與自由憲章第 33 條使立法機關得明示宣告法律不受憲章第 2 條或第 7 至 15 條的拘束。作者仍將加拿大計為強式，理由是該條款實際上甚少援用\pg{1356--1357}。註 34 回應 Goldsworthy「該條款足以回應民主疑慮」之說：援用該條款等於要立法機關表示自己不認為憲章權利重要，而法院與立法機關的對立，通常是雙方都重視該權利而對其內容理解不同\pg{1357}。

\qaentry{第三部之三}{作者接受的違憲審查}

\begin{qaquestion}
作者是否反對一切形式的違憲審查？
\end{qaquestion}

第參部分前段，作者表示尊重一項論證：議員未必看得出法案所涉的權利問題，也未必預見法律適用時會發生何種權利問題，因此宜有一套機制，讓公民在問題出現時提請眾人注意\pg{1370}。他接著寫：

\begin{quote}
But this is an argument for weak judicial review only, not for a strong form of the practice in which the abstract question of right that has been identified is settled in the way that a court deems appropriate.\pg{1370}
\end{quote}

這項論證所支持的，是英國那種由法院宣告存在重要權利問題的制度，或由檢察總長審查法律案並公開指出其中權利問題的更弱安排；註 63 引紐西蘭權利法案法第 7 條\pg{1370}。

\qaentry{第三部之四}{另三項區分}

\begin{qaquestion}
除強弱之外，第壹部分還作了哪些區分？作者如何取捨？
\end{qaquestion}

\begin{longtable}{@{}L{0.30\linewidth}L{0.62\linewidth}@{}}
\toprule
區分 & 作者的取捨 \\
\midrule
\endhead
以權利為對象的審查與一般的憲法審查（1357--1358） & 本文處理前者。聯邦主義事項由大法官投票決定，作者另表保留；註 35 指出最著名的違憲審查判決 Marbury v.\ Madison 與個人權利無關 \\
事後的個案審查與事前的抽象審查（1358--1359） & 作者表明專注於事後審查，對事前審查不再多述 \\
普通法院與專設的憲法法院（1359） & 以 Goodridge 為普通法院行使審查之例，對專設憲法法院只寫一段 \\
\bottomrule
\end{longtable}

\begin{readernote}
後兩項區分與中華民國的制度直接相關：憲法法庭是專設的憲法法院，國家最高機關與立法委員現有總額四分之一以上聲請的法規範憲法審查（憲法訴訟法第 47 條、第 49 條）不以具體案件為前提。本文第肆部分 A 節批評的「個案取向」，在這類審查中本不存在；可以請聽者評估，本文其餘的論證對這類制度的適用程度。
\end{readernote}

% 第四部：論證的前提（第貳部分，1359–1369）。出處：細目次第三節；W06_1019_導讀問答.md「第一項假設」一則。
\chapter{論證的前提：四項假設}

\qaentry{第四部之一}{四項假設的內容}

\begin{qaquestion}
本文的論證以哪四項假設為前提？作者為何要把論證寫成附條件的形式？
\end{qaquestion}

第貳部分前言列出四項假設：

\begin{quote}
We are to imagine a society with (1) democratic institutions in reasonably good working order, including a representative legislature elected on the basis of universal adult suffrage; (2) a set of judicial institutions, again in reasonably good order, set up on a nonrepresentative basis to hear individual lawsuits, settle disputes, and uphold the rule of law; (3) a commitment on the part of most members of the society and most of its officials to the idea of individual and minority rights; and (4) persisting, substantial, and good faith disagreement about rights (i.e., about what the commitment to rights actually amounts to and what its implications are) \ldots\pg{1360}
\end{quote}

接著寫：「I shall argue that, relative to these assumptions, the society in question ought to settle the disagreements about rights that its members have using its legislative institutions.」\pg{1360}摘要第 2 段已說明，論據並非絕對或無條件，以社會有運作良好的民主制度、多數公民認真對待權利為條件\pg{1346}；任一條件不成立，論證可能不成立，留待第柒部分處理。

\begin{longtable}{@{}L{0.16\linewidth}L{0.35\linewidth}L{0.22\linewidth}L{0.16\linewidth}@{}}
\toprule
假設 & 內容與限定 & 後文取用處 & 不成立時 \\
\midrule
\endhead
一、民主制度（1361--1362） & 成年人普選、定期公平選舉、大型審議機關、兩院與委員會、政黨；運作良好而未必完美；只指程序，不預設立法內容大體正義 & 第伍部分的立法機關答覆；第柒部分 & 1402--1403 \\
二、司法制度（1363--1364） & 獨立而運作良好的法院，已有權審查行政行為；非民選、非代議；回應具體訴訟、對抗式程序、遵循先例、合議庭以多數決作成決定 & 第肆部分 A 節；第伍部分的法院答覆 & 1402 \\
三、對權利的承諾（1364--1366） & 多數成員與公職人員認真對待個人與少數者的權利；權利懷疑論屬例外；已有一部權利法案，是否入憲、是否剛性不問 & 第陸、柒部分 & 1398；1404 \\
四、關於權利的分歧（1366--1369） & 就有哪些權利、其內容為何存在實質而善意的分歧，及於權利的核心適用 & 第參部分；第肆部分 B 節 & 1404 \\
\bottomrule
\end{longtable}

\qaentry{第四部之二}{第一、二項假設的平衡}

\begin{qaquestion}
第一項假設是否把立法程序理想化？作者如何避免乞題？
\end{qaquestion}

第一項假設限定其程度：

\begin{quote}
In general, I am assuming that the democratic institutions are in reasonably good order. They may not be perfect and there are probably ongoing debates as to how they might be improved.\pg{1361}
\end{quote}

同節末段把立法的內容排除在假設之外：作者不假定合理民主的立法機關所制定的法律大體良好或正義，他對立法與選舉安排運作良好的說法只涉及程序價值，不涉及結果價值\pg{1362}。提出第一項假設時，作者即以第二項假設平衡，設定法院同樣運作良好，表明不願乞題\pg{1362}。第肆部分前言另自承：對立法機關的樂觀描述是刻意的，用以對應法院支持者拒絕對法院作實際描述時的樂觀描述\pg{1379}。

\qaentry{第四部之三}{權利分歧與權利懷疑論}

\begin{qaquestion}
第四項假設所說的分歧，與道德相對主義或權利懷疑論有何不同？權利法案能否解決這種分歧？
\end{qaquestion}

作者所舉的分歧議題有墮胎、平權措施、重分配、犯罪嫌疑人的權利、宗教寬容、少數者的文化權利、選舉言論與支出的管制\pg{1367}。這些分歧出於善意，作者引 Rawls《政治自由主義》所稱的判斷的負擔\pg{1368}；承認分歧與道德相對主義無關\pg{1368}。分歧的雙方都認真對待權利，由第三項假設保證。權利法案與爭點有關，但不解決爭點；設有違憲審查時，分歧被構作成條文的解釋問題\pg{1368--1369}，第肆部分 B 節由此展開。

\begin{readernote}
第三、四兩項假設合起來，把「多數不尊重少數權利」的情形排在核心情形之外，留下的只有「雙方都尊重權利、對內容意見不同」的情形。違憲審查的辯護者最常舉的例子多半屬於前一種，聽者若覺得本文避開了真正的難題，可以指向這兩項假設與第柒部分。
\end{readernote}

% 第五部：評價決定程序的判準（第參部分，1369–1376）。出處：細目次第三節參；英文底稿 Quellen/waldron_source.tex。
\chapter{評價決定程序的判準：論證的形式}

\qaentry{第五部之一}{為何需要一個決定程序}

\begin{qaquestion}
在四項假設之下，社會為何必須另設一個決定程序？「立法有效，除非侵害權利」為何不能充當決定規則？
\end{qaquestion}

權利分歧的雙方都需要一個共同行動的基礎，權利問題也在立法場域出現\pg{1369--1370}。待答的問題是：立法機關就權利問題所為的決議應否具決定性，抑或由司法機關事後複核\pg{1371}。作者對決定程序提出兩項形式要求\pg{1371--1372}：

\begin{enumerate}
\item 「立法機關的決定有效，除非侵害權利」不能充當決定規則，因為社會成員爭執的正是某一法律案是否侵害權利。作者引 Hobbes 為據：決定程序的運作必須平息爭議，不得重燃當初使人需要決定程序的那些爭議\pg{1371}。
\item 社會成員雖對權利有分歧，仍須共享一套關於決定程序之正當性的理論\pg{1371}；程序的設計必須獨立於它要解決的那項分歧\pg{1373}。
\end{enumerate}

對「程序本身亦有分歧」的反對意見（Christiano 的無窮後退），作者答以：不論在決定程序上是否也有分歧，都必須處理正當性問題\pg{1371--1372}。任何程序都會出錯，Wollheim 弔詭（同一公民可以同時主張某政策不應施行、又因多數選擇而應施行）是政治理論的一般問題，違憲審查同樣會作出錯誤決定\pg{1372}。

\qaentry{第五部之二}{結果面理由與程序面理由}

\begin{qaquestion}
評價決定程序的理由分為哪兩類？各自的定義與限制為何？
\end{qaquestion}

\subsection{程序面的理由}

程序面的理由問的是何人有權參與決定，與決定的結果無涉。作者以父母與路人之例說明：孩子犯錯該不該管教，由父母決定，街上的路人或公車上的乘客無權決定，即使我們認為路人較可能作出更好的決定\pg{1372--1373}。政治平等與投票權屬於這一類。

\subsection{結果面的理由}

作者的定義是：

\begin{quote}
Outcome-related reasons, by contrast, are reasons for designing the decision-procedure in a way that will ensure the appropriate outcome (i.e., a good, just, or right decision).\pg{1373}
\end{quote}

結果面的理由不得以特定的爭議結果為目標，否則就是乞題。作者接受一種節制的版本：選擇最可能觸及權利真相的程序，不論真相為何\pg{1373}；他在此引 Kavanagh\pg{1373--1374}。

\subsection{兩類理由的規範性質}

Raz 以結果面理由為唯一值得考慮的理由，作者稱之為教條主義，不採\pg{1374}。作者認為兩類理由都帶有義務論上的急迫性；Dworkin 把參與的價值說成「參與上的後果」，作者認為這嚴重低估了參與權的觀念\pg{1375}。

\qaentry{第五部之三}{第參部分預告的結論}

\begin{qaquestion}
兩類理由權衡下來，作者預告的結論是什麼？
\end{qaquestion}

作者以兩種問法設定比較的問題，其一取自 Michelman（註 76），其二為作者將其前後倒置而成，接著寫：

\begin{quote}
What I will argue, in Part IV, is that the outcome-related reasons are at best inconclusive. They are important, but they do not (as is commonly thought) establish anything like a clear case for judicial review. The process-related reasons, however, are quite one-sided. They operate mainly to discredit judicial review while leaving legislative decisionmaking unscathed.\pg{1375}
\end{quote}

\begin{qafigure}{兩類理由對兩個機關的檢驗結果（1375--1393）。\label{fig:reasons}}
% 兩類理由 × 兩個機關（1375、1377–1393）。三格灰階，決定勝負的一格用強調色。
\begin{tikzpicture}[qa,cell/.style={nbox,text width=52mm,minimum height=17mm}]
  \node[hdr] at (3.15,1.25) {結果面的理由（第肆部分）};
  \node[hdr] at (8.85,1.25) {程序面的理由（第伍部分）};
  \node[hdr,anchor=east] at (0.15,0) {立法機關};
  \node[hdr,anchor=east] at (0.15,-2.0) {法院};
  \node[cell] at (3.15,0) {代表結構使受影響者的處境得以體察（1377--1378）；議會辯論可直接處理道德問題（1384--1385）};
  \node[cell] at (8.85,0) {公平選舉與多數決，給予平等的發言權與平等的決定權威（1387--1389）};
  \node[cell] at (3.15,-2.0) {個案取向、文本取向、說明理由三項優點，經檢驗不成立或兩方均有（1379--1386）};
  \node[cell,kbox] at (8.85,-2.0) {法官的民主資格遠不及議員；大法官之間的簡單多數決欠缺公平性的論證（1391--1392）};
  \node[lbl,anchor=north] at (3.15,-3.05) {至多未有定論};
  \node[lbl,anchor=north] at (8.85,-3.05) {一面倒，不利於違憲審查};
\end{tikzpicture}

\end{qafigure}

\begin{readernote}
這張表是全文的骨架：第肆部分填左欄，第伍部分填右欄。導讀時先放這張表，再分兩部講，聽者比較容易跟上。
\end{readernote}

% 第六部：結果面的理由（第肆部分，1376–1386）。出處：細目次第三節肆；W06_1019_導讀問答.md 案件表與理論表（已查核）；英文底稿。
\chapter{結果面的理由：法院是否較能得出正確的權利決定}

\qaentry{第六部之一}{第肆部分的出發點}

\begin{qaquestion}
第肆部分從誰的主張出發？作者為何說結果面的理由對兩個機關都有？
\end{qaquestion}

第肆部分從 Raz 的出發點與法律程序學派關於機構能力的論辯起論\pg{1376}。Raz 指控立法機關受宗派利益驅使；作者答以兩點。其一，同樣的宗派壓力也說明法院對權利的忽視，正文舉 Korematsu、Schenck、Dred Scott、Prigg 等案（註 82），並引通常堅定支持違憲審查的 Tribe 在九一一之後的評論：危機時期不宜對司法寄予過高的期望\pg{1377}。其二，立法機關若普遍如此，即與第三項假設不相容，屬非核心情形\pg{1377}。Raz 自己也承認結果面的理由有相反的方向：受決定影響的人較能體察處境，立法機關的代表結構正為培養這種體察\pg{1377--1378}。

民主程序本身已含有為結果而作的調整，例如選舉權的限制與立法機關的組成，屬 Rawls 所稱的不完全的程序正義\pg{1378}。作者接著指出，違憲審查的支持者對司法決定的描述不切實際\pg{1379}，以下逐一檢驗被歸給法院的三項優點。

\qaentry{第六部之二}{被歸給法院的三項優點}

\begin{qaquestion}
支持者認為法院在哪三方面較能得出正確的權利決定？\\
作者如何逐一回應？
\end{qaquestion}

\begin{longtable}{@{}L{0.16\linewidth}L{0.30\linewidth}L{0.46\linewidth}@{}}
\toprule
優點 & 支持者的主張 & 作者的回應 \\
\midrule
\endhead
A. 以個案為取向（1379--1380） & Moore：法官每天面對具體個人的處境，較能得到道德洞見 & 「But this is mostly a myth.」案件上訴到最高審級時，原權利人的痕跡幾已消失，論辯圍繞抽象的權利問題；原告常由倡議團體挑選\pg{1379--1380}。立法程序同樣考量個案，Megan's Law 即因知名個案而立法；「困難案件造就惡法」若屬實，立法機關較能評估個案與影響數百萬人的一般權利問題之間的關係\pg{1380} \\
B. 以權利法案的文本為取向（1380--1382） & 法院依附權利法案的條文推理 & 條文的措辭往往為迴避制定時的分歧而選；解釋理論的論辯排擠對道德問題的直接論辯；法院因自身正當性的疑慮而依附文本；文本所遺漏的權利（如社會經濟權利）同樣扭曲推理 \\
C. 說明理由（1382--1386） & 法院必須公開說明理由 & 議員同樣說明理由，載於 Hansard 與 Congressional Record；法院的理由以條文、先例為中心，成因是法院須證明自己有權作成該決定 \\
\bottomrule
\end{longtable}

\qaentry{第六部之三}{1966 年英國下議院墮胎法案辯論與 Roe 案}

\begin{qaquestion}
作者為何拿英國國會的一場辯論與 Roe v.\ Wade 判決對比？對比的內容為何？
\end{qaquestion}

作者寫到，在權利的善意分歧中真正的爭點常被推到邊緣：在 Roe v.\ Wade 五十頁的意見書中，處理生育權與隱私之關係的道德重要性只有幾段，處理胎兒地位的段落多半在呈現各種意見的分歧\pg{1383}；註中說該意見書含大量的法律史與社會史。

對照的是 1966 年英國下議院 Medical Termination of Pregnancy Bill 的二讀辯論。作者讀過議事錄，約一百頁的辯論直接處理墮胎的道德問題，連反對陣營的議員也稱許其立場受到尊重的傾聽\pg{1384--1385}；註 106 舉投票反對該法案的天主教議員 Norman St.\ John-Stevas，稱辯論的水準無愧於下議院最高的傳統。作者接著問，反對墮胎的一方何曾稱許最高法院在 Roe 案中對其立場的關注與尊重\pg{1385}。

\qaentry{第六部之四}{原則的論壇與常態、病態之分}

\begin{qaquestion}
作者如何回應 Dworkin「法院是原則的論壇」之說？\\
立法機關的缺陷為何不計入比較？
\end{qaquestion}

對 Dworkin 的回應，作者指出法院的推理以解釋與教義為主，英國的立法辯論能直接處理道德問題\pg{1385--1386}。他承認立法的推理也可能是 a disgrace：

\begin{quote}
\ldots so too can legislative reasoning be a disgrace, as legislative majorities act out of panic, recklessly, or simply parrot popular or sectarian slogans in their pseudo-debates.\pg{1386}
\end{quote}

他接著問兩個機關的缺陷何者屬常態、何者屬偏差。法院依先例、文本、教義推理正是被期待的行為方式，第二項假設所描述的就是依這些標準運作良好的法院；倉促或宗派性的立法則不屬立法機關被設立來做的事，不應假定於核心情形\pg{1386}。

\begin{readernote}
法院的缺陷計入常態，在 B、C 兩節拿來與立法機關比較；立法機關的缺陷劃為病態，移到第柒部分。這項劃分的依據是兩個機關設立的目的，作者未就兩者的實際運作作經驗上的比較。若以設立目的為準，法院的設立目的同樣可以說成就權利問題說理，法條主義也就可以歸為偏差。討論時可從這一段提問。
\end{readernote}

% 第七部：程序面的理由（第伍部分，1386–1395）。出處：英文底稿；W06_1019_導讀問答.md 理論表的調和論證一列（已查核）。
\chapter{程序面的理由：撇開結果，誰有權作終局決定}

\qaentry{第七部之一}{正當性理論須答覆的兩個問題}

\begin{qaquestion}
作者用虛擬的公民 C\textsubscript{n} 提出哪兩個問題？為何正當性理論必須答覆兩者？
\end{qaquestion}

政治決定通常由一組指定的人 \{C\textsubscript{1}, C\textsubscript{2}, \ldots, C\textsubscript{m}\} 以指定的程序作成。C\textsubscript{n} 不在其內，可以問：「Why them? Why not me?」即使接受了答案，C\textsubscript{n} 還可以問：在所用的程序中，為何不給與我看法相同的決定者更大的權重？正當性理論須同時說明成員的組成與所用的程序何以正當\pg{1387}。

\begin{qafigure}{C\textsubscript{n} 的兩個問題與兩個機關的答覆（1387--1392）。\label{fig:cn}}
% C_n 的兩個問題 × 兩個機關（1387–1392）。三格灰階，第二問對法院一格用強調色。
\begin{tikzpicture}[qa,cell/.style={nbox,text width=52mm,minimum height=15mm}]
  \node[hdr] at (3.15,1.2) {立法機關（約五百人）};
  \node[hdr] at (8.85,1.2) {聯邦最高法院（九人）};
  \node[hdr,anchor=east,align=right] at (0.15,0) {第一問\\[-1pt]{\scriptsize\mdseries 為何由這些人}};
  \node[hdr,anchor=east,align=right] at (0.15,-1.8) {第二問\\[-1pt]{\scriptsize\mdseries 為何採此程序}};
  \node[cell] at (3.15,0) {公平選舉：C\textsubscript{n} 與全體公民受到平等對待（1387）};
  \node[cell] at (8.85,0) {總統提名、參議院同意的間接民主資格，與議員相較遠為不足（1391）};
  \node[cell] at (3.15,-1.8) {多數決的公平性論證：中立、平等，給每一意見最大可能的權重（1388）};
  \node[cell,kbox] at (8.85,-1.8) {五票勝過四票；大法官不代表任何人，多數決在此欠缺道德基礎（1391--1392）};
\end{tikzpicture}

\end{qafigure}

\qaentry{第七部之二}{立法機關的答覆}

\begin{qaquestion}
對立法機關，C\textsubscript{n} 的兩個問題如何得到答覆？
\end{qaquestion}

作者把兩個問題寫成具體的形式：為何由大約五百位議員決定一個影響我與另外兩億五千萬人的權利問題？即使接受，為何不給與我看法相同的議員更大的權重\pg{1387}？第一問由公平選舉的理論答覆，C\textsubscript{n} 在決定誰得參與這類決定時，與全體公民受到平等對待；第二問由多數決（MD）的公平性論證答覆\pg{1387--1388}。作者寫道，MD 比其他規則更能在爭議的結果之間保持中立、平等對待參與者，並在給予所有意見同等權重的前提下，給每一意見最大可能的權重\pg{1388}。代表之間的多數決與選舉合起來，近似全體公民之間的多數決，給予平等的發言權與平等的決定權威\pg{1388--1389}。現實制度的缺陷由第一項假設處理，作者寫，正當性理論「will talk about reasonable fairness, not perfect fairness」\pg{1389}。

\qaentry{第七部之三}{法院的答覆與 Scalia 的引文}

\begin{qaquestion}
同樣兩個問題向聯邦最高法院提出時，答覆為何？作者為何長段引用 Scalia？
\end{qaquestion}

作者說正當性是比較性的，於是向法院提出 C\textsubscript{n} 的兩個問題\pg{1389}：為何由這九人決定此事？即使由他們決定，為何採用現行的程序，而不採給予與 C\textsubscript{n} 看法相同的大法官更大權重的程序\pg{1390}？作者寫，這兩問對法院遠比對議員難答，接著長段引用 Scalia 大法官在 Planned Parenthood v.\ Casey 案的不同意見：Scalia 對遊行、信件、抗議等「政治壓力」同感困擾，認為成因在於法院不再依文本與傳統實務，而依本院所稱的「說理的判斷」，那不過是哲學偏好與道德直覺\pg{1390--1391}。

第一問的答覆是大法官經總統提名、參議院同意而有間接的民主資格；作者認為與議員相較遠為不足\pg{1391}。第二問的處境更糟：

\begin{quote}
But in the end it comes down to head-counting: five votes defeat four in the U.S. Supreme Court, irrespective of the arguments that the Justices have concocted. \ldots{} MD is appropriate for persons who have a moral claim to insist on being regarded as equals in some decision-process. But I cannot see any moral basis for this claim in the case of Supreme Court Justices. They do not represent anybody. Their claim to participate is functional, not a matter of entitlement.\pg{1391--1392}
\end{quote}

文獻對法院內部的多數決少有討論；若以 Condorcet 陪審團定理為法院辯護，該辯護屬結果面理由，而依該定理，群體愈大，增進愈多，反而有利於人數較多的立法機關\pg{1392~及註~122}。

\qaentry{第七部之四}{五項調和論證及其反駁}

\begin{qaquestion}
辯護者提出哪些論證，主張違憲審查與民主可以調和？\\
作者如何逐一反駁？
\end{qaquestion}

\begin{longtable}{@{}L{0.42\linewidth}L{0.50\linewidth}@{}}
\toprule
調和論證 & 作者的反駁 \\
\midrule
\endhead
法官只執行有民主資格的權利法案 & 權利法案的抽象條文不能解決分歧\pg{1393} \\
預先承諾：社會像 Ulysses 把自己綁在桅桿上，法官是確保繩索未鬆的船員 & 社會並未承諾任何特定的權利理解；一旦出現另一種理解，既有的預先承諾為何仍應拘束並不清楚\pg{1393--1394}；註 123 引 Elster, Ulysses Unbound 對先前 Ulysses and the Sirens 的修正 \\
得修改權利法案而未修改，等於默示認可 & 修改通常需要特別多數；英國或紐西蘭式的法律位階權利法案也因其在政治文化中的地位而提高修改的負擔；起點由判決設定，不修改不能視為默示認可\pg{1394} \\
法官經提名同意而具有民主資格（Eisgruber） & 法官的民主資格遠不及議員；以民主資格接受法官，反而削弱把違憲審查當作一種獨特而有價值之決策形式的正面論證\pg{1394} \\
訴訟是公民參與的額外途徑 & 這種途徑不受政治平等原則的約束\pg{1395} \\
\bottomrule
\end{longtable}

\begin{readernote}
導讀時可以點出第四項的反駁：辯護者一旦以法官的民主資格為理由，就承認了程序面的比較，而在這項比較上法官必然輸給議員。作者據此說它削弱了違憲審查自己的正面論證。
\end{readernote}

% 第八部：多數暴政（第陸部分，1395–1401）。出處：英文底稿；譯本附錄圖「決定上的多數與議題上的多數」。
\chapter{多數暴政}

\qaentry{第八部之一}{「暴政」一詞為何不能直接支持違憲審查}

\begin{qaquestion}
「多數暴政」是反對立法至上最常見的理由。作者如何處理「暴政」一詞？由多數所加諸，是否使暴政加重？
\end{qaquestion}

作者先指出，任何權利分歧的雙方都可以指對方的主張構成暴政；任何決定程序都可能作成錯誤的決定，法院也是\pg{1395--1396}。由多數所加諸，並不使暴政加重。他的看法是：多數決的一面反而減輕了暴政，因為該決定至少有一件事不是暴政式的，它沒有以暴政的方式把某些人排除在平等參與之外\pg{1396}。

\qaentry{第八部之二}{決定上的多數與議題上的多數}

\begin{qaquestion}
作者區分哪兩組多數與少數？表決落敗，何時才構成多數暴政？
\end{qaquestion}

作者區分「決定上的」（decisional）多數、少數與「議題上的」（topical）多數、少數\pg{1397}。該處的註說明，用 topical 一詞，是因為這些人的權利與利益正是該決定的議題；此一用語是寬鬆的，但足以讓人看出，並非每個在權利議題上投給落敗一方的人，都是權利受該決定不利影響的群體成員\pg{1397}。

作者以自己為例：我參與表決，我投的一方落敗，我因此屬於決定上的少數。要證明我遭受了暴政，還須證明兩件事：(1) 該決定確屬錯誤，對受影響者的權利帶有暴政的含意；(2) 我屬於權利受此錯誤決定不利影響的議題上的少數\pg{1397}。即使我屬於議題上的少數，情況也不必然改變，因為人們包括議題上的少數，未必擁有他們自認擁有的權利\pg{1398}。

\begin{qafigure}{決定上的多數、少數與議題上的多數、少數（1397--1401）。\label{fig:majorities}}
% 決定上與議題上的多數、少數（1397–1398、1401）。三格灰階，兩種少數重合的一格用強調色。
\begin{tikzpicture}[qa,cell/.style={nbox,text width=52mm,minimum height=15mm}]
  \node[hdr,align=center] at (3.15,1.25) {議題上的少數\\[-1pt]{\scriptsize\mdseries 權利與利益為該決定之議題者}};
  \node[hdr] at (8.85,1.25) {其他人};
  \node[hdr,anchor=east,align=right] at (0.15,0) {決定上的多數\\[-1pt]{\scriptsize\mdseries 表決勝方}};
  \node[hdr,anchor=east,align=right] at (0.15,-1.8) {決定上的少數\\[-1pt]{\scriptsize\mdseries 表決敗方}};
  \node[cell] at (3.15,0) {議題上的少數也有人投給勝方；墮胎權、平權措施的支持與反對者跨越性別與族群（1401）};
  \node[cell] at (8.85,0) {勝方若主要由議題上的多數組成，對應左下格的重合（1401）};
  \node[cell,kbox] at (3.15,-1.8) {兩種少數重合：決定若確屬錯誤，構成多數暴政；普遍如此即屬非核心情形（1397--1398）};
  \node[cell] at (8.85,-1.8) {只是表決落敗：自己的意見未被採行，本身不構成暴政（1397--1398）};
\end{tikzpicture}

\end{qafigure}

作者的小結：

\begin{quote}
To sum up, tyranny of the majority is possible. But the term should not be used simply to mark the speaker's disagreement with the outcome of a majority decision. The most fruitful way of characterizing tyranny of the majority is to say that it happens when topical minorities are aligned with decisional minorities.\pg{1398}
\end{quote}

依此界定，多數暴政是非核心情形的特徵：人們除自己的權利之外，甚少關心少數或個人的權利。作者強調這與第三項假設不相容，不應同時談論一個承諾權利的社會、又說多數暴政在其中是普遍的可能\pg{1398}。

\qaentry{第八部之三}{對 Dworkin「背景條件」說的回應}

\begin{qaquestion}
Dworkin 主張尊重權利是任何決定制度具正當性的背景條件。\\
作者如何回應？
\end{qaquestion}

Dworkin 在 Freedom's Law 主張，人民的權利若未受尊重，民主決定本身即屬暴政，因為尊重權利是任何政治決定制度具正當性的背景條件\pg{1398--1399}。註中說明 Dworkin 自己表示這不是支持違憲審查的論據。作者提出兩點\pg{1400}。

\begin{enumerate}
\item Dworkin 似乎認為，決定若涉及民主本身，就不必追問作成決定的程序。作者不同意：若關於多數決程序的決定以排除女性參與的程序作成，以平等為據的反對不會因爭點是多數決的正當性而失去資格；即使爭點就是程序，我們仍在乎程序價值。
\item 「民主程序只在彼此尊重權利的人之間具正當性」有兩種讀法：須持有並實踐對彼此權利的正確見解，或認真對待彼此的權利並真誠設法確認其內容。前者過強，沒有任何可想像的政治制度滿足它；作者接受後者，而依後者，本文所設想的社會已滿足這項前提。
\end{enumerate}

\qaentry{第八部之四}{任何人不得為自己案件之法官}

\begin{qaquestion}
有人援引 nemo iudex in sua causa，主張權利問題不應由人民自己終局決定。\\
作者如何回應？
\end{qaquestion}

作者說，幾乎任何可想像的決定規則，最後都有人就自己的案件作決定；除非上訴的鏈條無窮無盡，總有某人或某機關的決定是終局的，而該終局者總是就自己見解的可接受性作裁斷\pg{1400--1401}。這項論證若有意義，是指議題上的多數不應同時是決定上起決定作用的多數；兩者重合時確有正當的疑慮，若普遍如此則屬非核心情形。作者以墮胎與平權措施為例，指出這類重合在美國違憲審查常處理的案件中甚少發生：許多女性支持墮胎權，許多男性也支持，也有許多女性反對；平權措施的支持與反對者同樣跨越族群\pg{1401}。

% 第九部：非核心情形與結論（第柒部分、結論，1401–1406）。出處：英文底稿；W06_1019_導讀問答.md「孤立而隔絕的少數」一則（已查核）。
\chapter{非核心情形與結論}

\qaentry{第九部之一}{假設不成立時論證的效力}

\begin{qaquestion}
四項假設不成立時，本文的論證如何？違憲審查是否因此就成立？
\end{qaquestion}

作者先答覆「本文的假設不切實際」之說：第三項假設相當容易滿足，因為違憲審查的支持者幾乎都預設該社會有一部與公民見解有實質關聯的權利法案；第一項假設並不要求完美\pg{1401--1402}。接著寫：

\begin{quote}
In cases in which the assumptions fail, the argument against judicial review presented in this Essay does not go through. As I emphasized in Part II, my argument is a conditional one. However, it does not follow that judicial review of legislation is defensible whenever the assumptions fail.\pg{1402}
\end{quote}

理由仍是比較性的：可能有不以這些假設為條件的其他反對論；違憲審查也可能無助於改善特定的處境；一個社會的司法若與其立法同樣腐敗、同樣帶有偏見，設立違憲審查並不適當\pg{1402}。

第一項假設不成立時，作者提出兩個問題：能否改善立法機關的處境？應否把重要權利問題的終局決定權交給法院\pg{1402--1403}？兩者並非全然獨立：法院作為後盾，可能使議員更難養成負責任的文化，作者自承 How far this is true, I don't know\pg{1403}。

\qaentry{第九部之二}{孤立而隔絕的少數}

\begin{qaquestion}
第柒部分為何以 Carolene Products 案腳註四收尾？作者如何界定 discrete 與 insular，又加了哪些限制？
\end{qaquestion}

1938 年 United States v.\ Carolene Products Co.\ 案腳註四寫：

\begin{quote}
[P]rejudice against discrete and insular minorities may be a special condition, which tends seriously to curtail the operation of those political processes ordinarily to be relied upon to protect minorities \ldots\pg{1403}
\end{quote}

作者稱這是描述違憲審查論據有若干說服力之非核心情形的極佳方式\pg{1403}，接著界定兩個形容詞\pg{1404}。

\begin{longtable}{@{}L{0.14\linewidth}L{0.42\linewidth}L{0.36\linewidth}@{}}
\toprule
用語 & 作者的界定 & 與第陸部分的關係 \\
\midrule
\endhead
discrete（孤立） & 獨立於政治決定而存在的少數，即議題上的少數 & 議題上的少數（1397） \\
insular（隔絕） & 成員與社群其他人少有共同利益，無從組成聯盟促進自身利益 & 兩種少數的重合（1398） \\
\bottomrule
\end{longtable}

對 Stone 所提的偏見要件，作者認為普遍的偏見與第三、四項假設不相容；重點在把各種敵意與對權利的合理分歧區分開來\pg{1404}。他另加三點限制：

\begin{enumerate}
\item 並非每個少數都適格，決定上的少數、議題上的少數、長期在表決中落敗者都不當然適格\pg{1403--1404}。
\item 要看司法上的多數是否同樣帶有偏見。支持違憲審查的論點假定同情少數權利的人集中於政治菁英，菁英中的法官較議員不受民意報復；同情若較多存在於一般人之間，選舉反而是把這份支持送進立法機關的管道\pg{1404--1405}。
\item 長期效果。美國 1950 與 1960 年代，法院攻擊種族隔離，最後起決定作用的是國會的民權立法，主要的分野在聯邦與州之間\pg{1405}。
\end{enumerate}

作者的結論是，不應把腳註四當作撬出一般違憲審查制度的槓桿\pg{1405--1406}。

\qaentry{第九部之三}{結論的兩項指控}

\begin{qaquestion}
結論把主張限定在什麼範圍？對違憲審查提出哪兩項指控？對支持者有什麼要求？
\end{qaquestion}

主張的範圍：

\begin{quote}
I have tried to show why rights-based judicial review is inappropriate for reasonably democratic societies whose main problem is not that their legislative institutions are dysfunctional but that their members disagree about rights.\pg{1406}
\end{quote}

結論第 2 段寫，一般的立法程序能以負責而審議的方式處理權利分歧，法院的終局審查在此之上所增加的，只是一種相當侮辱人的剝奪參政權（disenfranchisement），以及以法條主義遮蔽分歧中的道德問題（legalistic obfuscation）\pg{1406}。第 3 段承認，在特殊的病態、失靈的立法機關、腐敗的政治文化、種族主義與其他根深柢固偏見的遺緒之下，這些代價或許值得暫時承擔；但支持者應坦白以此為理由，並帶著 a degree of humility and shame，不應把違憲審查當作尊重權利的典範向外推廣\pg{1406}。

\begin{readernote}
結論兩項指控與導論的兩條戰線\pg{1353}一一對應：法律技術的遮蔽對應第肆部分，剝奪參政權對應第伍部分。導讀收尾時可以把導論與結論並排，讓聽者看到全文首尾的對應。
\end{readernote}

% 第十部：本課脈絡中的位置。出處：課程倉 notes/docs/1151/seminar/W02_0921.md、W05_1012_預習.md；
% 1151/seminar/readings/literature-mentions.tsv；各篇文字層 readings/W04、W12、W13/text/。
\chapter{本課脈絡中的位置}

\qaentry{第十部之一}{與前幾週指定文獻的關係}

\begin{qaquestion}
本文與第 2 週 Tremblay、第 4 週 Bickel 與 Fallon、\\
第 5 週 Kelsen 與 Schmitt 各有什麼關係？
\end{qaquestion}

\begin{qafigure}{本課指定文獻與本文的引用關係，依各篇文字層查得。\label{fig:course}}
% 本課指定文獻與本文的引用關係。灰階；本文用強調色。關係直接寫在各框第二行，不用斜放的線上標籤。
\begin{tikzpicture}[qa,wk/.style={nbox,text width=36mm,minimum height=10mm},
    bus/.style={figGray,line width=0.6pt}]
  \node[wk] (tr) at (0,2.55) {W02　Tremblay 2003\\{\scriptsize\color{dimink}未引本文；舉證責任的方向與本文相反（導讀人對照）}};
  \node[wk] (bi) at (0,0.85) {W04　Bickel 1962\\{\scriptsize\color{dimink}本文 1349 引其「對抗多數的難題」}};
  \node[wk] (fa) at (0,-0.85) {W04　Fallon 2018\\{\scriptsize\color{dimink}第 1 章註 25 引本文 1387--1389}};
  \node[wk] (ks) at (0,-2.55) {W05　Kelsen 1929、Schmitt 1931\\{\scriptsize\color{dimink}本文未引；均問誰作終局決定（導讀人對照）}};
  \node[kbox,text width=22mm,minimum height=10mm] (w) at (5.4,0) {W06\\Waldron 2006};
  \node[wk] (gh) at (10.75,0.85) {W12　Ghosh 2010\\{\scriptsize\color{dimink}以本文為主要對話對象之一}};
  \node[wk] (be) at (10.75,-0.85) {W13　Benhabib 2020\\{\scriptsize\color{dimink}討論 Waldron 的政治處境與預先承諾}};
  \coordinate (lb) at (2.95,0);
  \coordinate (rb) at (7.85,0);
  \foreach \n in {tr,bi,fa,ks} \draw[bus] (\n.east) -| (lb);
  \draw[bus] (lb) -- (w.west);
  \foreach \n in {gh,be} \draw[bus] (\n.west) -| (rb);
  \draw[bus] (rb) -- (w.east);
\end{tikzpicture}

\end{qafigure}

\subsection{第 2 週：Tremblay（2003）}

Tremblay 早於本文三年發表。第 2 週的課堂記錄整理其主張：違憲審查不需要特定授權，司法在社會中已經行使有效的強制力並具有政治權威，單是這個事實就構成承認其一般正當性的充分理由；特定正當性命題推定不正當，除非證明有特定授權，一般正當性命題則推定正當，除非能指出某個特定的道德原則使它不被允許。\footcite[525--562]{tremblay2003}

\begin{readernote}
本文第伍部分要求每一個決定程序都能向不同意的 C\textsubscript{n} 說明為何由這些人、以此程序決定\pg{1387}，等於讓違憲審查承擔說明的責任。兩文的舉證責任分配方向相反，可以請聽者比較：Tremblay 從「所有國家機關都受同一原則拘束」起算，Waldron 從「決定程序須有雙方均能接受的正當性理論」起算。
\end{readernote}

\subsection{第 4 週：Bickel 與 Fallon}

本文導論引 Bickel 的「對抗多數的難題」與「脫軌制度」作為既有反對論的代表\pg{1349~及註~7、8}。\footcite{bickel1962} Fallon 的專書第 1 章在註 25 引本文 1387 至 1389 頁，即 C\textsubscript{n} 的兩個問題與立法機關的答覆；\footcite[181 註 25]{fallon2018} Fallon 另於 2008 年以〈The Core of an Uneasy Case for Judicial Review〉為題回應本文。\footcite{fallon2008}

\subsection{第 5 週：Kelsen 與 Schmitt}

本文未引 Kelsen 與 Schmitt。第 5 週的 Schmitt《憲法的守護者》第一章指出，法學上真正的爭點是由誰決定（quis judicabit）；若以一個可裁判一切疑義的憲法法院取代憲法上有權的機關，這個國事法院就會是與國會、帝國總統及帝國政府並列的一個政治機關。\footcite[100--101]{vinx2015} Kelsen 1928 年的報告則從法秩序的位階構造推出憲法法院。\footcite[22--78]{vinx2015}

\begin{readernote}
Schmitt 與 Waldron 都追問終局決定者的身分，理由不同：Schmitt 以司法的界限為據，認為憲法爭議的決定本質上是政治的；Waldron 承認法院有能力審查立法\pg{1363}，問題在於它應否審查、其判斷應否終局，依據是政治平等。兩人主張的決定者也不同，一為帝國總統與政府，一為民選的立法機關。這一則對照屬導讀人所作，原文之間沒有引用。
\end{readernote}

\qaentry{第十部之二}{後續各週如何引用本文}

\begin{qaquestion}
本課後續各週的指定文獻中，哪些引用、回應本文？
\end{qaquestion}

依課程倉的互見計數，後續指定文獻提及 Waldron 之姓最多者為第 12 週的 Ghosh（2010）與第 13 週的 Benhabib（2020）。

\begin{itemize}
\item 第 12 週 Ghosh（課綱作 Lilian Castro，刊本作者為 Eric Ghosh）以本文為主要的對話對象之一：刊本 330 至 331 頁摘述本文的程序面與結果面理由，引本文 1372、1373 頁；其後回應本文的程序面論證與結果面論證。\footcite{ghosh2010}
\item 第 13 週 Benhabib 討論 Waldron 的「政治的處境」（circumstances of politics）、「立法的尊嚴」與「參與式多數決」，並轉述 Waldron 對預先承諾之說的反對。\footcite[509--513]{benhabib2020}
\end{itemize}

% 第十一部：2006 年後的發展與回應。出處：課程倉 W06_1019_導讀問答.md「本文發表後與本文論證相關的發展」一則；
% 原件在 1151/seminar/readings/W06/sources/ 與翻譯倉 JudicialReview/Quellen/raw/發展_20261007/。
\chapter{2006 年後的發展與回應}

\qaentry{第十一部之一}{作者的後續著作與論者的回應}

\begin{qaquestion}
本文發表後，作者本人與其他論者有哪些直接相關的著作？
\end{qaquestion}

本文第伍部分指出，法院內部的簡單多數決欠缺公平性的論證\pg{1391--1392}。作者 2014 年在 Yale Law Journal 第 123 卷發表〈Five to Four: Why Do Bare Majorities Rule on Courts?〉\nobreak，專以這個問題為題。\footcite{waldron2014}

Fallon 2008 年以〈The Core of an Uneasy Case for Judicial Review〉為題回應本文，\footcite{fallon2008} 第 4 週指定的 Fallon 專書第 1 章另在註 25 引本文 1387 至 1389 頁。\footcite[181 註 25]{fallon2018}

\qaentry{第十一部之二}{美國：本文所引兩案的後續}

\begin{qaquestion}
本文開篇的同性婚姻判決與多次引用的 Roe v.\ Wade，其後各有什麼發展？
\end{qaquestion}

本文以麻薩諸塞州 2003 年的 Goodridge 案開篇\pg{1348}。2015 年 6 月 26 日，聯邦最高法院在 Obergefell v.\ Hodges 以五比四判決，各州禁止同性婚姻或拒絕承認他州合法締結的同性婚姻，違反第十四修正案。\footcite{obergefell2015}

本文多次以 Roe v.\ Wade 為例：導論說支持選擇權的一方因擔心動搖 Roe，不願考慮反對違憲審查的論證\pg{1351}；第肆部分說 Roe 五十頁的意見書中處理墮胎道德問題的只有幾段\pg{1383}。2022 年 6 月 24 日，聯邦最高法院在 Dobbs v.\ Jackson Women's Health Organization 推翻 Roe 與 Casey，法院意見寫：

\begin{quote}
Roe and Casey must be overruled, and the authority to regulate abortion must be returned to the people and their elected representatives.\footcite{dobbs2022}
\end{quote}

\begin{readernote}
Dobbs 把墮胎的規範權交還各州的立法機關，與本文主張的決定場域相同；作成這項決定的仍是以多數決投票的大法官，可對照第伍部分對法院內部簡單多數決的批評\pg{1391--1392}。支持與反對 Roe 的兩方在 Dobbs 前後對違憲審查的態度，也可以拿 Charles Black 所說的 a sometime thing\pg{1351} 來對照。
\end{readernote}

\qaentry{第十一部之三}{紐西蘭與加拿大的變化}

\begin{qaquestion}
本文用來界定弱式與居間形式的紐西蘭、加拿大，其後制度與實務有何變化？
\end{qaquestion}

\subsection{紐西蘭}

本文寫作時，紐西蘭法院即使自行宣告法律與權利法案不相容，該宣告對立法程序也不生法律效果\pg{1356}。2018 年，紐西蘭最高法院在 Attorney-General v Taylor [2018] NZSC 104 確認高等法院有權作成不相容宣告，該案涉及 2010 年禁止受刑人投票的修正法。2022 年國會通過 New Zealand Bill of Rights (Declarations of Inconsistency) Amendment Act 2022，新增第 7A 條：上級法院（senior court，含高等法院、上訴法院與最高法院）就法律與權利法案不一致所為的宣告確定後，檢察總長應於眾議院第六個開議日以前提出通知，使眾議院知悉該宣告。\footnote{New Zealand Bill of Rights (Declarations of Inconsistency) Amendment Act 2022 (2022 No 45), \url{https://www.legislation.govt.nz/act/public/2022/45/en/latest}。}

\subsection{加拿大}

本文把加拿大計為強式審查，理由是儘管條款在實務上甚少援用\pg{1356--1357}。2019 年魁北克省為通過禁止部分公務員在職務上配戴宗教標誌的第 21 號法案，預先援用該條款；其後安大略省（2021、2022）、薩斯喀徹溫省（2023）等也相繼援用，涉及的事項包括第三方選舉廣告支出的限制（安大略 2021）、禁止加拿大公共雇員工會所屬教育工作者罷工（安大略 2022，旋即廢止）、未滿 16 歲學生在校更改名字與代名詞須得家長同意（薩斯喀徹溫 2023）等。\footnote{\url{https://www.lawnow.org/where-are-my-charter-rights/}；\url{https://www.cp24.com/news/canada/2025/09/17/alberta-defends-notwithstanding-clause-in-arguments-for-bill-21-supreme-court-case/}。}

\begin{readernote}
紐西蘭的不相容宣告自此對立法程序發生作用，在本文的分類上與英國的弱式審查距離縮短；加拿大歸類所依據的前提「甚少援用」，在 2019 年以後已有變化。
\end{readernote}

% 第十二部：討論題與中華民國的對照。出處：課程倉 1151/seminar/my/導讀_W06_Waldron_細目次與大綱.md 第四節「可在討論時提出的問題」；
% 盤點_Waldron概念之中華民國司法援用.md；W06_1019_導讀問答.md「孤立而隔絕的少數」一則。全部屬導讀人的提問。
\chapter{討論題與中華民國的對照}

以下各題均為導讀人的提問，依各部的順序排列；中華民國的對照為導讀人所作，本文未言及中華民國。

\qaentry{第十二部之一}{關於攻擊的標的}

\begin{qaquestion}
\begin{enumerate}
\item 第 1 週老師的分法以宣告違憲後法律的效力為準：宣告違憲即失效者為強式，宣告不符合憲法而效力交由政治部門處理者為弱式。本文以法院得否在個案中拒絕適用為準\pg{1354--1355}。兩種標準在英國、紐西蘭兩例上結論相同，在美國與中華民國的定期失效宣告上是否相同？
\item 本文對事前的抽象審查與專設憲法法院各只寫一段\pg{1358--1359}。第肆部分 A 節批評的個案取向在抽象審查中本不存在，其餘論證對憲法法庭的適用程度如何？
\end{enumerate}
\end{qaquestion}

\qaentry{第十二部之二}{關於四項假設}

\begin{qaquestion}
\begin{enumerate}
\item 第一項假設是否把立法程序理想化？作者把法院的法條主義算作常態、把倉促或宗派性的立法劃為病態\pg{1386}，這一劃分的依據是兩機關的設立目的，可否反過來適用於法院？
\item 中華民國就同性婚姻先後有司法院釋字第 748 號解釋與 2018 年的公民投票。以這段經過逐項檢驗四項假設，中華民國落在核心情形之內嗎？
\end{enumerate}
\end{qaquestion}

\qaentry{第十二部之三}{關於程序面的理由}

\begin{qaquestion}
\begin{enumerate}
\item 本文說法院的簡單多數決欠缺公平性的論證\pg{1391--1392}。114 年憲判字第 1 號以「集體議事的多數決原則」說明憲法訴訟法第 30 條第 1 項的評決門檻；湯德宗大法官在釋字第 770 號以「法院判決不備理由，即無正當性可言」為判決正當性的來源。兩者如何回應 Waldron？
\item Ulysses 之喻\pg{1393--1394}：制憲者或立法者所拘束者為後代，與 Ulysses 拘束自己者不同，跨世代的拘束以何為正當性的基礎？Elster《Ulysses Unbound》序言引 Jens Arup Seip：在政治上，人們從不試圖約束自己，只試圖約束他人。
\end{enumerate}
\end{qaquestion}

\qaentry{第十二部之四}{關於孤立而隔絕的少數}

\begin{qaquestion}
釋字第 748 號解釋理由書提高審查標準的理由之一：
\begin{quote}
同性性傾向者因人口結構因素，為社會上孤立隔絕之少數，並因受刻板印象之影響，久為政治上之弱勢，難期經由一般民主程序扭轉其法律上劣勢地位。
\end{quote}
湯德宗大法官在釋字第 692 號不同意見書則以該案被歧視者「要非所謂社會上『孤立而隔絕之少數』（discrete and insular minority）」為由，不提高審查基準。用本文第柒部分的三點限制\pg{1403--1405}檢驗兩者：同性性傾向者是否合於作者所界定的 discrete 與 insular？司法上的多數是否同樣受刻板印象影響？解釋作成之後的公民投票與立法經過，在作者所說的長期效果上意味什麼？
\end{qaquestion}

\qaentry{第十二部之五}{關於本文的適用範圍：釋字第 380 號}

\begin{qaquestion}
本文限定其對象為法院對立法的審查\pg{1353--1354}。釋字第 380 號審查的客體為教育部所訂大學法施行細則第 22 條第 1 項後段與第 3 項，理由為逾越大學法或未經大學法授權，不在本文設定的對象之內。
\begin{enumerate}
\item 本號理由書指明，共同必修科目「應本大學自治之精神由法律明文規定，或循大學課程自主之程序由各大學自行訂定」。立法院若以法律明定國文為共同必修，大學以講學自由爭之，在四項假設成立時即為本文所稱的核心情形。此時應由法院或立法機關作終局決定？
\item 董翔飛、施文森、曾華松三位大法官的不同意見書謂，立法者修訂大學法時如認共同必修損及學術自由或大學自治「當可明文予以禁止」，立法過程紀錄「並無是類討論，亦足證立法者無意排除上項規定」。本文所駁的第三項調和論證\pg{1394}同樣從立法者未表示反對推論其同意，兩者可否並列檢驗？
\end{enumerate}
\end{qaquestion}

\qaentry{第十二部之六}{關於中文學界的接受}

\begin{qaquestion}
中文文獻中，本文見於蘇彥圖 2018 與林建志 2014 的註；蘇彥圖 2017 的註引的是 Law and Disagreement 與 Constitutionalism: A Skeptical View，未引本文。湯德宗《權力分立新論》卷二增訂五版（2026，課綱指定）是否處理本文？可請教老師。
\end{qaquestion}


\backmatter
\printbibliography[heading=bibintoc,title={引用文獻}]

\end{document}
